\documentclass[aspectratio=169]{beamer}
\usepackage[utf8]{inputenc}
\usepackage[T1]{fontenc}
\usepackage{fontawesome}
\usepackage{subcaption}
\usepackage{graphicx}
\graphicspath{{assets/img/}}

\title{Applied Gitology}
\date{22 Fevrier 2019}
\author{Basile Pesin}

\usetheme{metropolis}

\begin{document}

\begin{frame}
\titlepage
\end{frame}

\begin{frame}
  \frametitle{Git, GitHub, kezako ?}
  \begin{center}
    \includegraphics[height=.25\paperheight]{git-logo.png}\hspace{.1\paperwidth}
    \includegraphics[height=.25\paperheight]{github-logo.png}
  \end{center}
  \begin{itemize}
  \item Git: Logiciel de gestion de versions (VCS) décentralisé. Gère des repositories (fichiers + historiques)
  \item GitHub: Hebergement et partage des repositories
  \end{itemize}
\end{frame}

\begin{frame}
  \frametitle{Repository}
  \begin{center}
      \includegraphics[height=.5\paperheight]{mes-repos.png}\\\vspace{.02\paperheight}
    Fichiers + historique
  \end{center}
\end{frame}

\begin{frame}
  \frametitle{Commit}
  \framesubtitle{Unité de changement}
  \begin{center}
    \includegraphics[width=.7\paperwidth]{commits.png}
  \end{center}
\end{frame}

\begin{frame}
  \frametitle{Commit}
  \framesubtitle{Ajout de fichier}
  \begin{center}
    \includegraphics[width=.8\paperwidth]{commit.png}
  \end{center}
\end{frame}

\begin{frame}
  \frametitle{Status}
  \framesubtitle{Quand on est perdu}
  \begin{center}
    \includegraphics[width=.8\paperwidth]{status.png}
  \end{center}
\end{frame}

\begin{frame}
  \frametitle{Gestion de branches}
  \begin{center}
    \includegraphics[height=.7\paperheight]{branches.png}
  \end{center}
\end{frame}

\begin{frame}
  \frametitle{Conflit et résolution}
  \begin{center}
    \includegraphics[width=.9\paperwidth]{conflict.png}\\
    Parfois on crée des conflits...\\\bigskip
    \includegraphics[height=.2\paperheight]{conflict_res.png}\\
    Mais on peut les régler !
  \end{center}
\end{frame}

\begin{frame}
  \frametitle{Remotes}
  \begin{center}
    \includegraphics[height=.7\paperheight]{push_pull.png}
  \end{center}
\end{frame}

\begin{frame}
  \frametitle{Remotes}
  \framesubtitle{Gestion sur GitHub}
  \begin{itemize}
  \item GitHub permet d'héberger les remotes
  \item Créer un repo ou en fork un autre
  \item Puis $git\:clone$
  \end{itemize}
\end{frame}

\begin{frame}
  \frametitle{Collaboration}
  \framesubtitle{GitHub et la collaboration}
  \centering
  GitHub est ``La plus grande communauté Open Source du monde'': Git et GitHub permettent aux développeurs de collaborer efficacement.\\
  \vspace{10pt}
  40 millions d'utilisateurs sur GitHub
\end{frame}

\begin{frame}
  \frametitle{Collaboration}
  \framesubtitle{En pratique (le GitHub flow)}
  \begin{enumerate}
  \item Trouver (ou ouvrir) une issue pour bug, nouvelle feature, etc...
  \item Forker le repo et créer une branche
  \item Coder !
  \item Creer une PR (typiquement 1 PR = 1 branche) pour discuter du travail en cours, puis le soumettre (review). Lire le CONTRIBUTING.md !!
  \item Si la PR est approuvée et mergée, cool ! Sinon (changes requested) -> retour à 3
  \end{enumerate}
\end{frame}

\begin{frame}
  \frametitle{GitHub Student Developer Pack}
  \centering
  \includegraphics[height=0.5\paperheight]{student-dev-pack.png}\\
  \url{https://education.github.com/pack}
\end{frame}

\begin{frame}
  \frametitle{GitKraken}
  \framesubtitle{Un GUI pour Git}
  \centering
  \includegraphics[height=0.75\paperheight]{gitkraken.png}
\end{frame}

\begin{frame}
  \frametitle{Plus d'informations}
  \begin{itemize}
    \item Des guides: \url{https://guides.github.com/}
    \item S'entrainer: \url{https://lab.github.com/}
    \item GitHub flow: \url{https://guides.github.com/introduction/flow/}
    \item La fiche aide mémoire distribuée
  \end{itemize}
\end{frame}

\end{document}
